\documentclass[aspectratio=169]{beamer}

\usetheme{Madrid}
\usecolortheme{default}
\setbeamertemplate{navigation symbols}{}

\usepackage{amsmath,amssymb}
\usepackage{booktabs}
\usepackage{graphicx}
\usepackage{listings}
\usepackage{xcolor}

\title[CIR+CVN]{CIR+CVN: Bridging LLM Semantic Understanding and Petri-Net Verification for Concurrent Programs}
\author{Anonymous Author(s)}
\institute{Anonymous Institution}
\date{\today}

\newcommand{\Cir}{\textsc{Cir}}
\newcommand{\Cvn}{\textsc{Cvn}}

\lstdefinelanguage{rust}{
  keywords={fn, let, mut, while, if, else, return, true, false, Arc, Mutex, Condvar},
  sensitive=true,
  morecomment=[l]{//},
  morecomment=[s]{/*}{*/},
  morestring=[b]",
}

\lstset{
  basicstyle=\ttfamily\tiny,
  keywordstyle=\bfseries,
  commentstyle=\color{gray},
  breaklines=true,
  showstringspaces=false,
  columns=flexible,
  keepspaces=true
}

\begin{document}

\begin{frame}
  \titlepage
\end{frame}

\begin{frame}{Problem Setting}
  \begin{itemize}
    \item Concurrency bugs arise from thread interleavings rather than a single local control path.
    \item Dynamic tools expose concrete failing schedules but cannot rule out unexplored interleavings.
    \item Source-level static verification must first recover shared-resource identity and protection relations.
    \item In realistic systems, that recovery is entangled with aliasing, ownership, heap indirection, and API idioms.
  \end{itemize}
\end{frame}

\begin{frame}{Why Source-Level Verification Is Hard}
  \begin{itemize}
    \item Lock-order analysis, protocol checking, and protection checking all depend on a reliable notion of resource identity.
    \item In Rust-like code, a verifier may need to track a mutex or condvar through allocation, cloning, borrowing, and closure capture.
    \item The alias-resolution step often dominates cost and precision before the verifier can even reason about interleavings.
  \end{itemize}
  \vspace{0.4em}
  \textbf{Paper stance:} we do \emph{not} claim to solve source-level alias analysis. We reformulate the problem so the formal back-end no longer depends on it.
\end{frame}

\begin{frame}{Key Idea: Model-First Verification}
  \begin{block}{Core reformulation}
    Ask an LLM to synthesize an explicit concurrency model from the specification, then apply formal verification to that model.
  \end{block}
  \begin{itemize}
    \item Step 1: generate an alias-free concurrency artifact with globally named resources.
    \item Step 2: statically validate the artifact.
    \item Step 3: translate it mechanically into a Petri-net-based verification model.
    \item Step 4: exhaustively explore bounded interleavings.
    \item Step 5: map counterexamples back to statement identifiers for targeted repair.
  \end{itemize}
\end{frame}

\begin{frame}{Two Core Formalisms}
  \begin{columns}[T]
    \begin{column}{0.48\textwidth}
      \begin{block}{\Cir}
        Concurrency Intermediate Representation
      \end{block}
      \begin{itemize}
        \item Statement-level concurrency model
        \item Alias-free, globally named resources
        \item Explicit protection mapping
        \item Stable statement IDs for diagnostics
      \end{itemize}
    \end{column}
    \begin{column}{0.48\textwidth}
      \begin{block}{\Cvn}
        Concurrency Verification Net
      \end{block}
      \begin{itemize}
        \item Weighted P/T Petri net
        \item Finite global variable store
        \item Three-valued guards for branching
        \item Exhaustive bounded state-space analysis
      \end{itemize}
    \end{column}
  \end{columns}
\end{frame}

\begin{frame}[fragile]{Running Example}
  \begin{columns}[T]
    \begin{column}{0.49\textwidth}
      \textbf{Rust bug}
\begin{lstlisting}[language=rust]
fn worker(shared: Arc<(Mutex<bool>, Condvar)>) {
  let (m, cv) = &*shared;
  let mut g = m.lock().unwrap();
  while !*g {
    g = cv.wait(g).unwrap();
  }
}

fn notifier(shared: Arc<(Mutex<bool>, Condvar)>) {
  let (m, cv) = &*shared;
  let mut g = m.lock().unwrap();
  cv.notify_one(); // bug
  *g = true;       // too late
}
\end{lstlisting}
    \end{column}
    \begin{column}{0.49\textwidth}
      \textbf{Generated \Cir}
\begin{lstlisting}
resources:
  m0:    { kind: Mutex }
  cv0:   { kind: Condvar, paired_with: m0 }
  ready: { kind: Var, init: false }
protection:
  ready: [m0]
worker:
  - { sid: w1, op: {lock: m0} }
  - { sid: w2, op: {read: ready} }
  - { sid: w3, op: {wait: {cv: cv0, mutex: m0}} }
notifier:
  - { sid: n1, op: {lock: m0} }
  - { sid: n2, op: {notify_one: cv0} }
  - { sid: n3, op: {write: {var: ready, val: true}} }
\end{lstlisting}
    \end{column}
  \end{columns}
\end{frame}

\begin{frame}{Pipeline Overview}
  \begin{center}
    \includegraphics[width=0.88\textwidth]{figures/graphviz.png}
  \end{center}
  \vspace{-0.4em}
  \begin{itemize}
    \item The loop is generate, statically check, translate, analyze, repair.
    \item Repairs are accepted only if they are bug-free \emph{and} preserve business goals.
  \end{itemize}
\end{frame}

\begin{frame}{Two-Layer Verification and Repair}
  \begin{table}
    \centering
    \footnotesize
    \begin{tabular}{@{}lll@{}}
      \toprule
      Bug Type & Detected By & Repair Strategy \\
      \midrule
      Missing drop & \Cir\ checker & Auto-fix \\
      Undefined resource & \Cir\ checker & LLM re-synthesis \\
      Spawn/join mismatch & \Cir\ checker & LLM re-synthesis \\
      Deadlock & \Cvn\ search & Counterexample-guided repair \\
      Signal loss & \Cvn\ search & Counterexample-guided repair \\
      Starvation & \Cvn\ search & Counterexample-guided repair \\
      \bottomrule
    \end{tabular}
  \end{table}
  \begin{itemize}
    \item Layer 1 filters malformed artifacts before formal analysis.
    \item Layer 2 targets deep interleaving bugs after the model is structurally valid.
  \end{itemize}
\end{frame}

\begin{frame}{Bug Detection and Goal-Aware Repair}
  \begin{itemize}
    \item Bug predicates include deadlock, signal loss, livelock, starvation, and channel block.
    \item Counterexamples are projected back to \Cir\ statement IDs through the anchor map $\mu$.
    \item The repair loop does not accept a model just because it is bug-free.
    \item It additionally checks whether all user-stated business goals are reachable in the bug-free \Cvn.
  \end{itemize}
  \begin{block}{Why goals matter}
    A repair can remove the bug by deleting required behavior. Goal reachability catches bug-free but semantically incomplete repairs.
  \end{block}
\end{frame}

\begin{frame}{Formal Results}
  \begin{itemize}
    \item Mechanical translation from \Cir\ to \Cvn.
    \item Forward simulation from \Cir\ executions to \Cvn\ executions.
    \item Backward simulation under concrete valuations.
    \item Deadlock correspondence and signal-loss precision results.
    \item Soundness result for the verification layer.
  \end{itemize}
  \vspace{0.4em}
  \textbf{Interpretation:} the theory supports the verification model, but the trust boundary remains at the generated \Cir, not arbitrary source code.
\end{frame}

\begin{frame}{Evaluation Setup}
  \begin{itemize}
    \item 9 bounded-concurrency patterns in total
    \item 6 buggy patterns, 3 bug-free baselines
    \item 5 LLMs across strong, frontier, and mini tiers
    \item Up to 5 rounds for generation and 5 rounds for repair
    \item Wall-clock evaluation on Apple M4 Pro / 24 GB RAM
  \end{itemize}
  \begin{block}{Research questions}
    RQ1: valid \Cir\ generation. \quad
    RQ2: detection, repair, and goal preservation. \quad
    RQ3: translation implementation correctness. \quad
    RQ4: source--\Cir\ trust boundary.
  \end{block}
\end{frame}

\begin{frame}{RQ1: \Cir\ Generation}
  \begin{table}
    \centering
    \footnotesize
    \begin{tabular}{@{}lcc@{}}
      \toprule
      Model & Success Rate & Avg. Rounds \\
      \midrule
      GPT-4o & 9/9 & 1.3 \\
      Claude 4 Sonnet & 9/9 & 1.1 \\
      DeepSeek-V3 & 9/9 & 1.6 \\
      Gemini 1.5 & 9/9 & 1.9 \\
      GPT-4o-mini & 8/9 & 2.4 \\
      \bottomrule
    \end{tabular}
  \end{table}
  \begin{itemize}
    \item Stronger models usually produce valid \Cir\ in 1--2 rounds.
    \item Complex multi-condvar and RwLock patterns are the hardest cases.
  \end{itemize}
\end{frame}

\begin{frame}{RQ2: Detection and Repair}
  \begin{columns}[T]
    \begin{column}{0.48\textwidth}
      \begin{block}{Detection}
        \begin{itemize}
          \item All 6 seeded bugs detected
          \item 3 baselines verified with no false positives
          \item Analysis finishes in under 20 ms
        \end{itemize}
      \end{block}
    \end{column}
    \begin{column}{0.48\textwidth}
      \begin{block}{Repair}
        \begin{itemize}
          \item GPT-4o, Claude 4: 6/6 bug-free repairs
          \item DeepSeek-V3, Gemini 1.5: 6/6 with more rounds
          \item GPT-4o-mini: 4/6
        \end{itemize}
      \end{block}
    \end{column}
  \end{columns}
  \vspace{0.4em}
  \begin{itemize}
    \item Frontier models need about 1.2--1.5 rounds on average.
    \item Intermediate regressions are mostly shallow structural errors and are caught by Layer 1.
  \end{itemize}
\end{frame}

\begin{frame}{Goal Reachability Matters}
  \begin{itemize}
    \item The paper explicitly separates \emph{bug-freedom} from \emph{semantic completeness}.
    \item Some repairs pass static checking and bug detection, yet silently remove required behavior.
    \item Goal reachability is the final acceptance gate that prevents these false successes.
  \end{itemize}
  \begin{block}{Key message}
    The most distinctive systems contribution is not only \Cir\ + \Cvn, but the goal-aware acceptance criterion in the repair loop.
  \end{block}
\end{frame}

\begin{frame}{Limitations and Trust Boundary}
  \begin{itemize}
    \item The verified object is the generated \Cir, not the original source program.
    \item The approach assumes bounded concurrency and finite-state abstraction.
    \item The evaluation uses representative patterns rather than large real-world systems.
    \item Source-to-\Cir\ conformance remains future work.
  \end{itemize}
\end{frame}

\begin{frame}{Takeaways}
  \begin{itemize}
    \item The paper proposes a model-first formulation of concurrency verification.
    \item LLMs are used for intent recovery and repair, not for exhaustive verification.
    \item \Cir\ removes aliasing from the formal layer by construction.
    \item \Cvn\ enables explicit state-space analysis over bounded interleavings.
    \item Goal reachability is essential to reject bug-free but semantically incomplete repairs.
  \end{itemize}
\end{frame}

\end{document}
